\label{sec:zones-as-community}

The United States has, by Census Bureau count, approximately 90,000 governmental
units~\citep{census2017govts}.
These range from the federal government and fifty states at the top to more than
35,000 special districts---school districts, fire districts, water districts, utility
districts, mosquito abatement districts, hospital districts---at the bottom.
For census tracts, the most relevant administrative zones are those that directly
shape the daily civic and fiscal experience of residents.
We consider six zone types, described below in order of their legal and fiscal
significance.

\subsection{School Districts}

School districts are the single most politically salient local administrative boundary
for families with children.
The school district determines which school a child attends, the quality of that
education, and the property tax levy that funds it.
In the 2019--2020 school year, local property taxes provided approximately 36\% of
all public school revenue in the United States, with the exact fraction varying
sharply by state and district~\citep{nces2022}.
School district bond elections mobilize voter participation at rates comparable to
presidential primaries in many suburban jurisdictions.
Parents organize, advocate, and vote as members of their school district community.

From a redistricting perspective, placing tracts from the same school district in
different legislative districts fractures the school's political constituency.
A school principal seeking support for a bond measure or capital improvement project
can no longer turn to a single legislative representative; instead, parents must
navigate multiple elected officials with different priorities and different committee
assignments.
This fragmentation is a recognized harm, and expert witnesses in redistricting
proceedings across California, Colorado, Virginia, and North Carolina have cited
it as evidence of inadequate communities of interest preservation.

\subsection{County Subdivisions}

County subdivisions---the Census Bureau's term for minor civil divisions, which include
towns, townships, boroughs, and municipalities---are the primary property tax-levying
unit in New England and the Midwest.
In the six New England states (Connecticut, Maine, Massachusetts, New Hampshire,
Rhode Island, Vermont) and in several Midwestern states (Michigan, Minnesota, Wisconsin,
Indiana), the town or township is the governmental unit that levies the property tax,
maintains local roads, operates the local cemetery, and in some cases provides fire
and police protection.
The county in these states is a secondary administrative unit with limited powers.

The ``town meeting''---the annual gathering at which residents of a New England town
vote directly on the municipal budget, zoning bylaws, and capital expenditures---is
the oldest continuously operating democratic institution in the United States.
A census tract placed in the same congressional district as the rest of its town
can participate collectively in state-level advocacy about the issues the town
meeting addresses; a census tract split from its town cannot.

In states where the county is the primary unit (most of the South and West),
county subdivision co-membership is still meaningful as a marker of municipal
service delivery, but it is less politically central than in New England.

\subsection{Electric Utility Service Territories}

Every square meter of the United States is served by an electric utility, which
is required to report its service territory to the Energy Information Administration
on Form 861~\citep{eia861}.
EIA Form 861 data, supplemented by the Homeland Infrastructure Foundation-Level
Data (HIFLD) electric utility service territory shapefiles, provides 100\% national
coverage at no cost.

Rate-payers in the same electric utility service territory are a recognized economic
community under state public utility commission regulation.
When a utility seeks a rate increase, it must hold public hearings and notify customers
in its service territory; those customers have standing to intervene in the regulatory
proceeding.
When a utility proposes to build a new power plant or transmission line, it is the
rate-payers in its service territory who bear the cost and have the right to comment.
The public utility commission's service territory is, in this sense, a legally
recognized community of shared economic interest.

Electric utility co-membership is also the most uniformly available zone type:
unlike school districts (which are relevant primarily to households with children)
or fire districts (which vary in coverage by state), electric utility co-membership
affects every resident and every business.
It is available for every census tract in the United States, making it the most
reliable zone dimension in the M.6 formula.

\subsection{Fire Districts}

Fire protection districts, or fire protection authorities, are special districts that
provide fire suppression, emergency medical services, and hazardous materials response.
They are funded primarily by property taxes levied within the district and, in some
states, by subscription fees.
TIGER/Line Special District files identify fire districts with the code 21.

Coverage of fire district data in TIGER/Line is approximately 60\% of states.
In rural areas, fire protection is frequently provided by volunteer companies organized
as special districts; in urban areas, fire protection is typically a city or county
function and does not appear as a separate district.
The graceful degradation property of the zone co-membership formula ensures that
the absence of fire district data in a state merely reduces $|Z|$ by one, without
distorting the score for the other available zone types.

\subsection{Water and Sewer Districts}

Water supply and wastewater treatment districts are among the most numerous special
districts in the United States.
They are identified in TIGER/Line Special District files with codes 22 through 29
(various water and wastewater authority types).
Coverage is approximately 40\% of states; in many jurisdictions, water service is
provided by a municipality or county rather than a special district, and therefore
does not appear in the TIGER/Line special district layer.

Water co-membership is a meaningful community signal where it exists.
Neighbors who share a water system share a stake in water quality, infrastructure
replacement schedules, and rate decisions made by the water authority's governing
board.
In regions with water scarcity---the Southwest, California's Central Valley---water
rights and water district governance are among the most contested local political
issues.

\subsection{Police Precincts}

Police precincts are the most locally variable zone type.
In cities with municipal police departments, precincts can be obtained from city and
county GIS portals, but there is no national data source.
In rural areas, law enforcement is typically a county function (the sheriff's
department), and the county serves as a fallback unit.

Because of the highly variable data availability, police precincts are the
lowest-priority zone type in the M.6 formula.
The algorithm includes precinct co-membership when data is available and substitutes
county co-membership as a fallback.
Given the approximately 30\% coverage of municipal precinct data, police precincts
will frequently not contribute to the zone score, and the formula degrades gracefully
to the other five zone types.
